% !TEX root = ../main.tex

この部では、Lean を使って仕様書の核となる性質を証明する方法を扱います。
実務の仕様書には、曖昧な表現、運用上の判断、外部システムへの依存、パフォーマンス上の制約、例外的な業務ルールなどが含まれます。
これらを最初から完全に形式化しようとすると、作業量が急速に増え、実務上の価値が見えにくくなります。

本書では、仕様書の中から「小さく、重要で、壊れると困る性質」を切り出します。
たとえば、口座残高が負にならないこと、送金前後で総残高が保存されること、割引計算後の金額が期待した範囲に収まること、特定のリファクタリングで計算結果が変わらないことなどです。
このような性質は、Leanの型、関数、述語、定理として表現しやすいものです。

まず、自然言語仕様を、入力、出力、状態、操作、事前条件、事後条件、不変条件に分けます。
次に、証明したい性質を選びます。
影響が大きく、小さな命題として切り出せる核を証明対象に選びます。
証明対象が大きすぎると、Leanのコードも仕様書も読みにくくなります。
逆に、証明対象を適切に絞ると、仕様の曖昧な部分が明確になり、設計レビューにも使いやすくなります。

この部の中心例は、口座操作と料金計算です。
口座操作では、残高、不変条件、入金、出金、送金を扱います。
料金計算では、税、割引、手数料、丸めなど、順序によって結果が変わる可能性のある処理を扱います。
これらは実務に近く、かつLeanで小さなモデルを作りやすい題材です。

この部を読み終えるころには、仕様書の中から証明可能な性質を抽出し、Leanで小さな検証モデルを作る手順を理解できます。
また、Lean の検証モデルと実装・運用上の検査を対応付けられるようになります。

\section*{この部で扱うこと}

\begin{itemize}
\item 自然言語仕様の分解
\item 型、関数、述語、定理への変換
\item 事前条件、事後条件、不変条件
\item 口座操作の仕様証明
\item 料金計算の仕様証明
\item テストと証明の役割分担
\item 検証モデルと実装・運用上の検査の対応
\end{itemize}
